\section*{Generative AI use}
I used Claude (Anthropic) and OpenAI Codex as assistants for this assignment.

\textbf{Prompts.}
\begin{itemize}
  \item I asked Claude to read the PA1 handout and notebook, complete the required
        implementation cells, run the notebook on the full preset, and draft the Task 1
        responses from the numbered outputs.
  \item I asked Claude to build a small Autoformer solution for the leaderboard task,
        including preprocessing, chronological validation, optional-data ablation,
        multi-seed runs, plots, and the final 168-value submission string.
  \item I asked Codex to review Claude's work against the assignment requirements,
        check for over-engineering or non-student-like writing, verify the leaderboard
        submission files, and make the final report submission-ready.
\end{itemize}

\textbf{What the assistant produced.}
\begin{itemize}
  \item Task 1: the three notebook implementations, the full-preset run, and first drafts
        of Responses 1A--4 from the exported Outputs.
  \item Task 2: the \texttt{pa1q2} package (data handling, from-scratch Autoformer,
        training with resumable early stopping, baselines, aggregation, figures,
        submission writer), the experiment plan and runs, and a first draft of this
        section.
  \item Final review: checks for missing fields, exact prediction count, declared
        parameter and epoch values, and consistency between the report and saved outputs.
\end{itemize}

\textbf{My edits and checks.}
\begin{itemize}
  \item I checked the notebook outputs, kept the full-preset figures and tables, and
        edited the responses to focus on the required comparisons instead of only
        describing the code.
  \item For Task 2, I used my own validation split to choose the final model, kept the
        optional-data ablation, and did not use the public leaderboard as a validation
        signal.
  \item I verified that \texttt{submission/predictions.txt} contains exactly 168 values
        in chronological order and that the declared parameter count and epoch count
        match the saved final run.
\end{itemize}
